% Sources: docs/protocol.md (v1.0 as executed; sections 1-6 frozen, scripts/protocol_freeze.py), src/trace_vr/arms.py,
% scripts/train_arm.py (constants LR, WARMUP, R, TARGETS), src/trace_vr/model.py (BASE_ID,
% BASE_REVISION), src/trace_vr/parse_solve.py, src/trace_vr/token_budget.py,
% src/trace_vr/bootstrap.py, docs/exposure_pilot_train.md (generated by scripts/exposure.py).
% tests/test_paper_setup.py re-checks every number below against those files.
% manifests/pilot.json, manifests/pilot_test.json (Scale paragraph).
% No model result is stated in this section.
\section{Setup}
\label{sec:setup}

\paragraph{Hypothesis.}
H1: training a small VLM only to output minimal repair edits on drawn directed graphs
improves its accuracy on \emph{non-edit} relational questions (reachability, critical edge)
about unseen root graphs more than fine-tuning the same base on those questions directly,
or on structural graph tasks, at equal samples and images. H0: any gain is explained by
generic fine-tuning, extra data, output length, resolution or evaluation shortcuts.
Reachability appears inside the repair goals, so the claim concerns transfer from an edit
objective to questions, not the learning of a new concept.

\paragraph{Base model and adaptation.}
All arms start from \texttt{Qwen/Qwen3.5-2B} at a pinned revision (\texttt{15852e8c}). Trained
arms use LoRA with rank $r=16$, $\alpha=16$, no dropout, on the attention, linear-attention
and MLP projections of the language model; the vision encoder is frozen. Optimisation:
AdamW, learning rate $2\times10^{-4}$, weight decay $0.01$, linear warm-up over $3\%$ of the
updates then linear decay, gradient clipping at $1.0$, batch size 1 with 4-step gradient
accumulation, loss on the answer tokens only. Every image is $512\times512$ pixels and costs
256 image tokens for every arm and split. Training seeds are 3407, 2027 and 9173.

\paragraph{Arms.}
\begin{description}
\item[A, frozen base.] No training; same parser as every other arm. With the prompt every
  other arm sees, the base reasons in prose and runs out of tokens, so its prompt was chosen on
  development data by a rule fixed beforehand (highest strict non-edit macro accuracy at the
  variant's own token limit, ties to the common prompt). On 20 development roots (80 non-edit
  questions) the common prompt scored $0.1375$ at 768 new tokens (lenient parse $0.6125$,
  truncation up to $0.90$ per task); one added sentence (``Do not explain your reasoning.
  Reply with the JSON object only.'') scored $0.6000$ after the question and $0.6250$ as a
  system message, both at 96 new tokens with no truncation. Arm A therefore uses the system
  message at 96 new tokens. The two short variants differ by two questions; the choice
  follows the rule and is not evidence that one is better. A is a reported floor, not a
  compared arm.
\item[B, QA-SFT.] One non-edit twin pair per root graph, alternating reachability and
  critical-edge questions.
\item[C, structural SFT.] An adaptation inspired by MCDGraph~\cite{mcdgraph} (not a
  reproduction): on exactly the images D sees, a fixed cycle of graph description (list every
  node and edge), masked infilling (name the withheld third of a listed edge set) and
  one-image graph matching (is this edge list exactly the drawn graph? negatives differ by one
  edit). The edge lists in C's prompts are its declared extra training input.
\item[D, TRACE.] Minimal repair only: the root image and its repair twin of every root
  graph, answer \texttt{edited} with one edit, \texttt{unchanged} or \texttt{impossible}.
\item[E, parse and solve.] An extraction adapter, \texttt{C\_desc} (the same base with LoRA,
  trained per seed only on the graph-description task, list every node and edge, on the images
  D sees), transcribes each image once; fixed cleaning rules, decided before any E output existed, turn it into a graph on which the exact
  solver answers every question. E is scored against the true graph like every other arm.
\end{description}
D2 (D with an auxiliary ``does the goal hold after this edit?'' task) is the first of the two
revisions the protocol allows after a failed pilot. Because removal edits make that task close
to the critical-edge question, D2 is not a pure transfer test and is reported separately.

\paragraph{Matched exposure.}
Per epoch on the pilot training split every trained arm sees the same 1{,}024 samples from
the same 512 root graphs and 1{,}024 distinct images, and all arms train for the same number
of epochs. Text length is not equalised; it is reported (Table~\ref{tab:exposure}). C and D2
use exactly D's images; B uses the same root graphs but other twins and shares 546 of D's
1{,}024 images.

\begin{table}[h]
\centering
\small
\begin{tabular}{lrrrrr}
\toprule
arm & samples & root graphs & images & instruction tokens & target tokens \\
\midrule
B  & 1{,}024 & 512 & 1{,}024 &  73{,}216 &  5{,}120 \\
C  & 1{,}024 & 512 & 1{,}024 &  92{,}146 & 28{,}592 \\
D  & 1{,}024 & 512 & 1{,}024 & 195{,}542 & 16{,}384 \\
D2 & 1{,}024 & 512 & 1{,}024 & 148{,}410 & 10{,}764 \\
\bottomrule
\end{tabular}
\caption{Training exposure per epoch on the pilot training split (generated by
\texttt{scripts/exposure.py}; text tokens from the base tokenizer, without the chat template).
Image tokens are $256$ per sample for every arm.}
\label{tab:exposure}
\end{table}

\paragraph{Scale.}
This is a pilot-scale study. Every trained adapter saw 512 training root graphs with 5 to 10
nodes for 3 epochs, and every choice was made on 200 development roots. The protocol planned a
main build of about 5{,}000 training roots and locked splits of 500, 500, 500 and 250 roots;
that build was not run. The locked test has 200, 200, 200 and 100 root graphs (ID, style-OOD,
structure-OOD, combined-OOD); its structure-OOD and combined-OOD graphs have 11 to 16 nodes.
Every statement in this paper is a statement about adapters of this size.

\paragraph{Decoding and scoring.}
Decoding is greedy. The primary score is strict: after trimming whitespace and at most one
Markdown code fence the response must be exactly one JSON object with the declared keys;
prose or a reasoning preamble is a format failure and counts as wrong. A lenient score (last
balanced JSON object) is reported next to it for every arm, never instead of it. The token
limit per arm is fixed by a rule set before measurement: one greedy development run at 768 new
tokens gives every shorter limit as a prefix, and the arm uses the smallest of 96, 192, 384
and 768 at which no task exceeds $2\%$ truncation.

\paragraph{Statistic and interpretation rules.}
The primary statistic is the unweighted mean of six accuracies, \{reachability, critical
edge\} $\times$ \{ID test, style-OOD test, structure-OOD test\}, over root and edited images.
Intervals come from a paired bootstrap over root graphs (2{,}000 resamples, fixed seed). Fixed
in advance: a clear positive requires D to exceed the stronger of B and C (chosen on
development data) by about 5 points, not driven by one style or seed; a drop of more than 2
points in either task mean rules out ``general improvement''; if E beats D this is reported.
Test labels are locked in code and opened once, after every choice above is frozen on
development data.
